\chapter{Cơ sở lý thuyết}
\label{chap:background}

Chương này trình bày nền tảng lý thuyết của tấn công giải ẩn danh: từ hạn chế của các phương pháp ẩn danh truyền thống, mô hình hình thức hóa bài toán, thuật toán Scoreboard-RH và phân tích lý thuyết về lượng tri thức cần thiết, đến tình huống điển hình Netflix Prize và mối liên hệ với an ninh di động. Nội dung chương dựa chủ yếu trên công trình gốc của Narayanan và Shmatikov~\cite{narayanan}, có bổ sung từ các công trình liên quan.

\section{Bài toán ẩn danh hóa micro-data}

\textit{Micro-data} là dữ liệu ở mức từng cá nhân: mỗi bản ghi mô tả một người và gồm nhiều thuộc tính. Khác với dữ liệu thống kê tổng hợp (chỉ công bố các đại lượng gộp như trung bình, tổng), micro-data giữ lại bản ghi riêng của từng người ngay cả sau khi đã ``ẩn danh hóa''. Các tập như lịch sử mua hàng, hồ sơ y tế, log tìm kiếm, dấu vết di chuyển GPS hay lịch sử xem phim đều là micro-data.

Nhu cầu công bố micro-data ngày càng lớn: phục vụ nghiên cứu y học, huấn luyện hệ khuyến nghị, minh bạch hóa dữ liệu công (``open data''). Để bảo vệ người dùng, đơn vị công bố thường \textit{loại bỏ thông tin định danh trực tiếp} (Personally Identifiable Information - PII) như họ tên, số căn cước, địa chỉ. Câu hỏi trung tâm của chương này là: \textit{liệu việc gỡ bỏ PII có thực sự bảo vệ được quyền riêng tư hay không?}

Câu trả lời là không. Hai đặc trưng của micro-data hành vi - \textit{số chiều cao} (high dimensionality) và \textit{độ thưa} (sparsity) - khiến mỗi bản ghi gần như là một ``dấu vân tay'' duy nhất, có thể bị liên kết ngược về danh tính thật chỉ với một lượng nhỏ thông tin bên lề.

\section{Các phương pháp ẩn danh truyền thống và hạn chế}

\subsection{k-anonymity}
Mô hình bảo vệ kinh điển nhất là k-anonymity~\cite{sweeney}. Đơn vị công bố chia thuộc tính thành hai nhóm:
\renewcommand{\labelitemi}{$-$}
\begin{itemize}
	\item \textit{Quasi-identifier} (QID): các thuộc tính tuy không định danh trực tiếp nhưng có thể liên kết với nguồn dữ liệu khác để truy ra danh tính (ví dụ: tuổi, giới tính, mã vùng bưu chính).
	\item \textit{Thuộc tính nhạy cảm}: thông tin cần bảo vệ (ví dụ: bệnh án, mức lương).
\end{itemize}
Dữ liệu được \textit{khái quát hóa} (generalization) và \textit{nén} (suppression) sao cho mỗi tổ hợp giá trị QID xuất hiện ở ít nhất $k$ bản ghi. Như vậy, một bản ghi không thể bị phân biệt khỏi $k-1$ bản ghi khác cùng nhóm QID.

Động lực của k-anonymity đến từ một vụ tấn công thực tế: năm 1997, L. Sweeney liên kết tập dữ liệu xuất viện ``ẩn danh'' của bang Massachusetts (gồm mã vùng, ngày sinh, giới tính) với danh sách cử tri công khai, và truy ra được \textit{bệnh án của chính thống đốc bang}~\cite{sweeney}. Chỉ ba thuộc tính tưởng như vô hại đã đủ định danh phần lớn dân số Hoa Kỳ.

\subsection{l-diversity, t-closeness và giới hạn chung}
k-anonymity bảo vệ khỏi việc \textit{định danh bản ghi}, nhưng không bảo vệ giá trị nhạy cảm. Nếu cả $k$ bản ghi trong một nhóm QID đều có cùng giá trị nhạy cảm (tấn công đồng nhất - homogeneity attack), kẻ tấn công vẫn suy ra được giá trị đó. Các cải tiến ra đời để vá lỗ hổng này:
\begin{itemize}
	\item \textbf{l-diversity}~\cite{ldiversity}: yêu cầu mỗi nhóm QID chứa ít nhất $l$ giá trị nhạy cảm ``đủ khác biệt''.
	\item \textbf{t-closeness}~\cite{tcloseness}: yêu cầu phân phối thuộc tính nhạy cảm trong mỗi nhóm gần với phân phối toàn cục (khoảng cách $\le t$).
\end{itemize}

Tuy nhiên, toàn bộ họ phương pháp này có một giả định nền tảng bị phá vỡ khi áp dụng cho micro-data hành vi: chúng giả định có thể \textit{tách trước} tập QID khỏi thuộc tính nhạy cảm. Với dữ liệu như lịch sử xem phim, giả định này sai về bản chất:
\renewcommand{\labelitemi}{$-$}
\begin{itemize}
	\item \textbf{Mọi thuộc tính đều có thể là QID.} Kẻ tấn công có thể biết \textit{bất kỳ} bộ phim nào về nạn nhân; không thể chỉ định trước một tập QID cố định.
	\item \textbf{Thất bại trên dữ liệu nhiều chiều.} Khi số thuộc tính lên tới hàng nghìn, để mọi tổ hợp xuất hiện $\ge k$ lần phải khái quát hóa mạnh tới mức dữ liệu mất hết giá trị sử dụng~\cite{narayanan}.
\end{itemize}
Chính vì vậy, tập Netflix Prize - như sẽ thấy - hoàn toàn \textit{không} $k$-ẩn danh với bất kỳ $k > 1$ nào.

\subsection{Các tiền lệ rò rỉ thực tế}
Ngoài vụ Massachusetts, hai vụ việc nổi tiếng khác minh họa rủi ro:
\begin{itemize}
	\item \textbf{AOL (2006):} AOL công bố 20 triệu truy vấn tìm kiếm ``ẩn danh'' (chỉ thay tên người dùng bằng số). Phóng viên New York Times đã truy ra danh tính người dùng số 4417749 chỉ từ nội dung các truy vấn của bà~\cite{aol}.
	\item \textbf{MovieLens (2006):} Frankowski và cộng sự liên kết các bài đăng công khai trên diễn đàn MovieLens với hồ sơ đánh giá nội bộ, cho thấy ngay cả ``đề cập công khai'' về phim cũng đủ gây rò rỉ~\cite{frankowski}.
\end{itemize}

\section{Độ thưa và độ tương tự}

Quan sát then chốt của Narayanan-Shmatikov là các tập dữ liệu hành vi thực tế đều thưa: mỗi người chỉ tương tác với một phần rất nhỏ trong không gian thuộc tính khổng lồ. Trên tập Netflix, mỗi thuê bao chỉ đánh giá vài chục đến vài trăm phim trong số hàng chục nghìn phim. Hệ quả là mỗi bản ghi gần như duy nhất: rất khó tìm được hai người có cùng tổ hợp phim và điểm. Chính sự duy nhất này là mầm mống của rò rỉ riêng tư.

Gọi $\mathrm{supp}(r)$ là \textit{support} của bản ghi $r$ - tập các thuộc tính khác rỗng (các phim mà người đó đã đánh giá). Phân phối số người đánh giá mỗi phim tuân theo quy luật ``đuôi dài'' (long-tail) hay luật lũy thừa: một số ít phim bom tấn được rất nhiều người xem, còn phần lớn phim chỉ có ít người xem. Các phim hiếm ở phần đuôi mang giá trị định danh cao.

Để đo mức ``giống nhau'' giữa hai bản ghi, ta dùng hàm tương tự $\mathrm{Sim}(r_1, r_2) \in [0,1]$. Một cách hình thức, bài báo định nghĩa tính thưa như sau:

\vspace{0.2cm}
\noindent\textbf{Định nghĩa (Độ thưa).} \textit{Một cơ sở dữ liệu $D$ được gọi là $(\epsilon, \delta)$-thưa theo độ tương tự $\mathrm{Sim}$ nếu}
\begin{equation}
	\Pr_{r}\big[\, \exists\, r' \neq r : \mathrm{Sim}(r, r') > \epsilon \,\big] \le \delta
\end{equation}
\textit{nghĩa là xác suất một bản ghi ngẫu nhiên có ``hàng xóm gần'' với độ tương tự vượt $\epsilon$ là không quá $\delta$.}\footnote{Bài báo gốc viết điều kiện trong ngoặc dưới dạng $\mathrm{Sim}(r,r')>\epsilon\ \forall r'\neq r$. Ở đây chúng em dùng lượng từ $\exists$ cho đúng với ngữ nghĩa mà chính bài báo mô tả - ``không tồn tại bản ghi nào khác đủ giống'' - tức xác suất một bản ghi có ít nhất một hàng xóm gần là rất nhỏ.}

\vspace{0.2cm}
Bài báo chứng minh tập Netflix thưa áp đảo: với đại đa số bản ghi, không tồn tại bản ghi nào khác có độ tương tự trên $0{,}5$ trong toàn bộ 500.000 bản ghi - ngay cả khi chỉ xét tập phim đã xem mà bỏ qua điểm số và ngày.

\section{Mô hình tấn công và định nghĩa rò rỉ riêng tư}

\subsection{Tri thức bổ trợ}
Kẻ tấn công có một tập \textit{aux} $= \mathrm{Aux}(r)$ - thông tin bên lề về nạn nhân $r$, gồm giá trị của một tập con nhỏ các thuộc tính. Điểm mấu chốt làm nên tên gọi ``robust'': aux không cần chính xác - điểm có thể sai lệch, ngày có thể chỉ biết gần đúng, thậm chí một vài mục hoàn toàn sai. Một cách hình thức, $\mathrm{Aux}$ được mô hình hóa như một hàm xác suất ánh xạ bản ghi thật sang phiên bản nhiễu mà kẻ tấn công nắm giữ. Kẻ tấn công có thể thu thập aux từ trò chuyện đời thường, mạng xã hội, hay đánh giá công khai của nạn nhân trên một dịch vụ khác.

\subsection{Hai kịch bản rò rỉ}
Bài báo phân biệt hai dạng vi phạm riêng tư:
\renewcommand{\labelitemi}{$-$}
\begin{itemize}
	\item \textbf{Giải ẩn danh ``đoán tốt nhất'' (best-guess):} thuật toán đưa ra \textit{một} bản ghi $r'$ và khẳng định đó là nạn nhân. Vi phạm xảy ra nếu $\mathrm{Sim}(r, r') \ge \theta$ với xác suất cao.
	\item \textbf{Giải ẩn danh ``danh sách nghi phạm'' (lineup):} thuật toán đưa ra một tập nhỏ ứng viên kèm phân phối xác suất, tương tự khái niệm ẩn danh trong mạng trộn (mix network). Mức rò rỉ được đo bằng \textit{entropy} của phân phối: kẻ tấn công cần thêm bao nhiêu bit thông tin để xác định duy nhất nạn nhân.
\end{itemize}
Vì dữ liệu thưa, bản ghi thỏa $\mathrm{Sim}(r, r') \ge \theta$ gần như chắc chắn là duy nhất - nên danh sách nghi phạm thường thu về đúng một người. Khi đó kẻ tấn công học được \textit{toàn bộ} thuộc tính của bản ghi đó, kể cả các thuộc tính nhạy cảm mà ban đầu hắn không biết.

\section{Thuật toán Scoreboard và Scoreboard-RH}

Thuật toán giải ẩn danh được mô tả như một \textit{khuôn mẫu} (template) gồm ba thành phần:
\begin{enumerate}[1.]
	\item \textbf{Hàm cho điểm} $\mathrm{Score}(aux, r')$: gán cho mỗi bản ghi ứng viên một điểm thể hiện mức khớp với aux.
	\item \textbf{Tiêu chí khớp} (matching): từ tập điểm, quyết định có một khớp hay không.
	\item \textbf{Chọn bản ghi} (selection): trả về bản ghi ``đoán tốt nhất'' hoặc phân phối xác suất.
\end{enumerate}

Thuật toán cơ sở Scoreboard gán điểm bằng độ tương tự nhỏ nhất giữa ứng viên và aux, rồi chọn nếu điểm vượt một hằng số $\alpha$. Cách này không bền vững: chỉ cần một thuộc tính trong aux hoàn toàn sai là điểm sụp đổ.

Phiên bản bền vững Scoreboard-RH (Robust, Histogram) cải tiến hàm cho điểm theo hai heuristic quan trọng:

\subsection{Trọng số theo độ hiếm}
Các thuộc tính hiếm được gán trọng số lớn hơn. Biết nạn nhân xem một bộ phim nghệ thuật ít người biết có giá trị định danh cao hơn nhiều so với biết họ xem một bom tấn. Trọng số định nghĩa:
\begin{equation}
	wt(i) = \frac{1}{\log |\mathrm{supp}(i)|}
	\label{eq:weight}
\end{equation}
với $|\mathrm{supp}(i)|$ là số người đã đánh giá phim $i$. Phim càng ít người xem, $wt(i)$ càng lớn.

\subsection{Cộng gộp tương tự rating và thời gian}
Với mỗi phim $i$ trong aux, điểm tích lũy theo cả độ lệch điểm số và độ lệch thời gian, dùng hàm mũ suy giảm:
\begin{equation}
	\mathrm{Score}(aux, r') = \sum_{i \in \mathrm{supp}(aux)} wt(i)\left( e^{-\frac{|\rho_i - \rho'_i|}{\rho_0}} + e^{-\frac{|d_i - d'_i|}{d_0}} \right)
	\label{eq:score}
\end{equation}
trong đó $\rho_i, d_i$ là điểm và thời điểm trong aux; $\rho'_i, d'_i$ là của ứng viên $r'$; $\rho_0, d_0$ là hằng số co giãn (bài báo chọn $\rho_0 = 1{,}5$ và $d_0 = 30$ ngày). Hàm mũ khiến các khớp ``gần đúng'' vẫn được tính điểm dương - đây là chìa khóa của tính bền vững: một phim vẫn đóng góp ngay cả khi ngày xem sai lệch nhiều.

\subsection{Tiêu chí eccentricity}
Để quyết định ứng viên điểm cao nhất có thực sự là nạn nhân hay chỉ là trùng hợp, thuật toán dùng eccentricity: khoảng cách giữa điểm cao nhất và điểm cao nhì, chuẩn hóa theo độ lệch chuẩn của toàn bộ phân phối điểm:
\begin{equation}
	\mathrm{ecc} = \frac{\max_1 - \max_2}{\sigma}
	\label{eq:ecc}
\end{equation}
Nếu $\mathrm{ecc} > \phi$ (ngưỡng, bài báo chọn $\phi = 1{,}5$) thì ứng viên cao nhất ``nổi bật'' hẳn so với đám đông và được kết luận là nạn nhân; ngược lại, thuật toán \textit{từ chối định danh} để tránh dương tính giả. Đây là cơ chế tự kiểm soát sai số quan trọng: nó cho phép thuật toán \textit{biết khi nào nó không biết}, nhờ đó mỗi lần khẳng định đều có độ tin cậy cao.

\section{Phân tích: cần bao nhiêu tri thức bổ trợ?}

Một đóng góp lý thuyết quan trọng của bài báo là định lượng tri thức bổ trợ tối thiểu để tấn công thành công. Kết quả chính có thể phát biểu trực giác như sau: với một cơ sở dữ liệu $N$ bản ghi đủ thưa, kẻ tấn công chỉ cần biết khoảng
\begin{equation}
	m \ \ge\ \frac{\log N}{-\log(1 - \epsilon)}
\end{equation}
thuộc tính của nạn nhân (mỗi thuộc tính khớp tới độ chính xác $1-\epsilon$) là đủ để định danh duy nhất bản ghi với xác suất cao. Điều đáng chú ý là $m$ chỉ tăng theo logarit của kích thước cơ sở dữ liệu: dù dữ liệu lớn gấp nghìn lần, lượng tri thức cần thiết hầu như không đổi. Với dữ liệu \textit{thưa}, $m$ còn nhỏ hơn nữa - đây là lý do chỉ cần 5-8 thuộc tính là đủ để giải ẩn danh hầu hết bản ghi trong các tập thực tế.

Bài báo cũng chứng minh tính bền vững một cách hình thức: ngay cả khi dữ liệu công bố đã bị thêm nhiễu (sanitization) và một phần aux bị sai, thuật toán vẫn không tạo ra khớp giả với xác suất cao, miễn là dữ liệu đủ thưa. Trực giác là nhiễu nhẹ không thể ``xóa nhòa'' tính duy nhất vốn rất mạnh của các bản ghi thưa.

\section{Tình huống điển hình: tập Netflix Prize}

Năm 2006, Netflix công bố tập dữ liệu gồm hơn 100 triệu lượt đánh giá của khoảng 500.000 thuê bao (đã loại bỏ thông tin định danh) để tổ chức cuộc thi trị giá 1 triệu USD cải tiến hệ khuyến nghị. Netflix khẳng định dữ liệu an toàn vì ``mọi thông tin định danh đã được gỡ bỏ''. Tập công bố thực chất chỉ là một mẫu nhỏ (dưới $1/10$) của cơ sở dữ liệu gốc.

\subsection{Kết quả giải ẩn danh}
Áp dụng Scoreboard-RH, Narayanan-Shmatikov chứng minh điều ngược lại. Các kết quả chính (sẽ được chúng em tái hiện ở Chương~\ref{chap:exp} trên MovieLens):
\renewcommand{\labelitemi}{$-$}
\begin{itemize}
	\item Biết 8 phim (trong đó 2 phim có thể sai) cùng ngày xem với sai số 14 ngày là đủ định danh duy nhất 99\% thuê bao.
	\item Chỉ với 2 phim cùng ngày (sai số 3 ngày), 68\% bản ghi đã bị định danh.
	\item Ngay cả khi không biết ngày, 84\% thuê bao vẫn bị định danh nếu kẻ tấn công biết 6 trong 8 phim nằm ngoài top 500 phim phổ biến nhất.
\end{itemize}

\subsection{Tấn công thực tế bằng dữ liệu IMDb}
Để chứng minh tính khả thi, các tác giả lấy tri thức bổ trợ từ đánh giá công khai trên IMDb: nhiều người dùng đánh giá phim công khai trên IMDb (thường dưới tên thật) đồng thời là thuê bao Netflix. Chỉ với một mẫu nhỏ khoảng 50 người dùng IMDb, thuật toán định danh được bản ghi Netflix của họ với eccentricity lên tới 28 và 15 độ lệch chuẩn - những khớp cực kỳ chắc chắn, gần như không thể là dương tính giả.

Hậu quả về riêng tư là nghiêm trọng: từ lịch sử Netflix bị phơi bày, có thể suy ra các thông tin \textit{nhạy cảm không hề công khai trên IMDb} - như quan điểm chính trị (qua các phim tài liệu đã xem), tôn giáo, hay xu hướng tính dục của nạn nhân. Đây không phải vi phạm giả định: chính những suy luận này về sau dẫn tới một vụ kiện tập thể chống Netflix~\cite{netflixsuit} (viện dẫn Đạo luật Bảo vệ Riêng tư Video - VPPA~\cite{vppa}) và khiến Netflix hủy bỏ cuộc thi Netflix Prize thứ hai.

\subsection{Tính riêng tư xuôi (forward privacy)}
Một hệ quả tinh tế: một khi một mẩu dữ liệu bị liên kết với danh tính thật, mọi liên kết \textit{về sau} đều dễ dàng hơn - kể cả khi nạn nhân tạo một định danh ảo hoàn toàn mới. Điều này tương tự khái niệm ``forward secrecy'' trong mật mã, nhưng ở chiều ngược lại và bất lợi cho người dùng: dữ liệu hành vi một khi đã rò rỉ thì không thể ``thu hồi''.

\section{Liên hệ với an ninh di động}

Bài học của Narayanan-Shmatikov không giới hạn ở dữ liệu phim, mà đặc biệt thời sự với an ninh trên thiết bị và ứng dụng di động - nơi dữ liệu hành vi người dùng được thu thập với mật độ dày đặc và cũng mang đặc tính thưa, nhiều chiều:
\renewcommand{\labelitemi}{$-$}
\begin{itemize}
	\item \textbf{Dấu vết di chuyển (mobility traces).} de Montjoye và cộng sự~\cite{demontjoye} chứng minh rằng chỉ cần 4 điểm thời gian-không gian (ví dụ vị trí gần đúng tại 4 thời điểm) là đủ để định danh duy nhất 95\% người dùng trong một tập 1,5 triệu người - cùng bản chất ``thưa nên duy nhất'' như tập Netflix. Vị trí được điện thoại di động ghi lại liên tục, khiến đây là rủi ro trực tiếp.
	\item \textbf{Siêu dữ liệu giao dịch.} Cũng nhóm này~\cite{demontjoye2} cho thấy 4 giao dịch thẻ tín dụng là đủ định danh 90\% người dùng - dữ liệu thường được các ứng dụng ví điện tử thu thập.
	\item \textbf{Vân tay thiết bị và cảm biến.} Tập hợp các thuộc tính tưởng vô hại của một thiết bị di động (danh sách ứng dụng đã cài, phiên bản hệ điều hành, đặc trưng cảm biến gia tốc, danh sách font, độ phân giải) cũng tạo thành một ``dấu vân tay'' thưa và nhiều chiều, cho phép theo dõi xuyên ứng dụng mà không cần định danh trực tiếp.
\end{itemize}

Trong mọi trường hợp, cơ chế tấn công là chung: \textit{tính thưa + số chiều cao $\Rightarrow$ tính duy nhất $\Rightarrow$ khả năng liên kết}. Vì vậy, việc nghiên cứu kỹ tấn công Scoreboard-RH trên một ví dụ trực quan (phim) cung cấp khung tư duy để hiểu và phòng vệ trước cả lớp rủi ro giải ẩn danh trên thiết bị di động.
